\section{Exceptional-width returns with the actual finite-block costs}
\label{sec:exceptional-partial-return}

The local converse of MA-PR-004 can be assembled globally after
an adaptive regularization.  For an unregularized test there is
an exact additional topological criterion.  Both statements use
the original adapted block norms and the actual Hilbert quotient.
No polynomial bound for a boundary selector, a reciprocal edge
distance, or an individual interpolation matrix is assumed.

Retain (MA-PR.1), with arbitrary \(u<v\), \(a\in\mathbb R\)
and \(s\ge0\), but do not assume (MA-PR.2).
The multiplicities \(r_\rho\) are exactly (MA-WP.3).
Write \(\nu_n=2+|n|\).  The full smooth blocks \(J_n\),
their holomorphic quotient norms \(\|\cdot\|_n\), the rectangles
\(\Omega_n\), and the entire lifts \(\mathcal T_n\) are those
of MT-M18, without changing their norms.
In particular
\[
 \begin{aligned}
 Q_\infty&\simeq s(J),&
 q_k((b_n))&=\sup_n\nu_n^k\|b_n\|_n,\\
 p_{N,j}(\mathcal T_n b)&\le C_{N,j}\nu_n^{D_{N,j}}\|b\|_n .
 \end{aligned}
 \tag{MA-PR.44}
\]
Let \(\mathfrak R_n:Q_\infty\to J_n\) be the full block
restriction.  The actual Hilbert block and its algebraic
coordinate space are
\[
 V_n=P_nQ_H,\qquad
 B_n=\prod_{\substack{\rho\in\Omega_n\\r_\rho>0}}
                      \mathbb C[z]/(z^{r_\rho}),\qquad
 \pi_n:J_n\longrightarrow B_n .
 \tag{MA-PR.45}
\]
Here \(P_n\) is the bounded convolution projection of
(MA-WP.11)--(MA-WP.12).  It need not be orthogonal.
The finite algebra \(B_n\) in (MA-PR.45) is not assigned an
independent Hilbert norm in what follows.

The continuous retained-jet map on the \emph{whole} Hilbert
quotient is denoted by \(\mathscr R_n:Q_H\to B_n\).
For \(v=\pi h\), its coordinates are
\[
 (\mathscr R_nv)_{\rho,j}
    =\frac1{j!}\int_{\mathbb R}x^j e^{(\rho-a)x}h(x)\,dx,
       \qquad 0\le j<r_\rho .
 \tag{MA-PR.46}
\]
The integral is absolutely convergent by (MA-WP.17).
These functionals annihilate \(J_H\), so (MA-PR.46) is
independent of the representative \(h\).
The exact block description in MA-WP gives
\[
 \mathscr R_nP_n=\mathscr R_n,\qquad
 \mathscr R_n|_{V_n}:V_n\xrightarrow{\ \sim\ }B_n,\qquad
 \mathscr R_ni=\pi_n\mathfrak R_n .
 \tag{MA-PR.47}
\]
All three are identities of actual continuous maps.

For \(g\in\mathcal E\), put \(G=\mathcal Mg\) and
\[
 b_n(g)=R_n(T_{-a}g)\in J_n,\qquad
 b_n(g)(w)=G(w-a)\quad\hbox{as a full Taylor class}.
 \tag{MA-PR.48}
\]
Call \(g\) \emph{eligible} if at every proper nonzero
boundary truncation it satisfies
\[
 G^{(j)}(\rho-a)=0
       \quad\bigl(0\le j<\min(m_\rho-r_\rho,r_\rho)\bigr),
       \qquad 0<r_\rho<m_\rho .
 \tag{MA-PR.49}
\]
This defines a closed linear subspace \(\mathcal E_{\rm el}\)
of \(\mathcal E\).  No condition is imposed at a factor
with \(r_\rho=0\), or with \(r_\rho=m_\rho\).
Whether any particular boundary or exterior factor occurs
is not asserted.

\paragraph{MA-PR-006 (an exact global lifting criterion).}
For eligible \(g\), let
\[
 \begin{aligned}
 \mathcal F_n(g)
   =\{L\in\operatorname{Hom}_{\mathbb C[w]}(B_n,J_n):
                        \pi_nL=M_{\pi_nb_n(g)}\},\\
 c_n(g)=\min_{L\in\mathcal F_n(g)}
                \|L\mathscr R_n\|_{\mathcal L(Q_H,J_n)} .
 \end{aligned}
 \tag{MA-PR.50}
\]
For \(B_n=0\), set \(c_n(g)=0\).
Every nonempty minimum in (MA-PR.50) is attained.
These are finite-dimensional quantities with the actual source
norm; they do not discard the norm of an oblique projection.

There is a continuous linear map \(S_g:Q_H\to Q_\infty\)
with
\[
 iS_g=C_g^{Q_H},\qquad
 S_gU_t^{Q_H}=e^{-at}U_t^{Q_\infty}S_g
                  \quad(t\in\mathbb R)
 \tag{MA-PR.51}
\]
if and only if \(g\) is eligible and
\[
                  \sup_n\nu_n^k c_n(g)<\infty
                         \quad\hbox{for every integer }k\ge0 .
 \tag{MA-PR.52}
\]
Every continuous map satisfying (MA-PR.51) is nuclear of
every positive order and has a continuous linear lift
from \(Q_H\) to \(\mathcal E\).
There is no claim of uniqueness at a truncated boundary factor.

The costs in (MA-PR.50) admit an explicit Gram description.
Use the Taylor coordinates of (MA-PR.46) and their ordinary
Euclidean scalar product solely for this formula.  Then
\[
 \begin{aligned}
 K_n&=\mathscr R_n\mathscr R_n^*>0,\\
 (K_n)_{(\rho,j),(\sigma,k)}
   &=\int_{\mathbb R}
       \frac{x^{j+k}e^{(\rho-a+\overline{\sigma}-a)x}}
              {j!\,k!\,w_{u,v,s}(x)}\,dx,\\
 c_n(g)&=\min_{L\in\mathcal F_n(g)}
             \|LK_n^{1/2}\|_{\mathcal L(\ell^2,J_n)} .
 \end{aligned}
 \tag{MA-PR.53}
\]
The target in the last line retains its original
holomorphic quotient norm.

\begin{proof}
\emph{The finite module problem and the actual norms.}
At a primary factor, write
\(A=\mathbb C[z]/(z^m)\) and \(B=\mathbb C[z]/(z^r)\).
A module map \(L:B\to A\) is determined by \(c=L(1)\)
and is well-defined exactly when
\[
        c\in\operatorname{Ann}_A(z^r)=z^{m-r}A .
 \tag{MA-PR.54}
\]
Its reduction is multiplication by \(G(\rho-a+z)\)
if and only if \(c\equiv G(\rho-a+z)\pmod {z^r}\).
This is solvable exactly under (MA-PR.49), as follows
by reducing \(z^{m-r}A\) modulo \(z^r\).
For \(r=m\) the map is forced to be multiplication by
the full class; for \(r=0\) there is no source factor.
Module maps between primary factors at distinct \(\rho\)
are zero, by the coprimality of their annihilating
polynomials.  Thus (MA-PR.49) is precisely the
nonemptiness of every \(\mathcal F_n(g)\).

The maps \(\mathscr R_n\) are surjective because their
restriction to \(V_n\) is an isomorphism.  Consequently
\(K_n\) is positive definite when \(B_n\ne0\).
The same Gram matrix can be computed on \(H\):
the quotient map \(\pi:H\to Q_H\) is a coisometry and
the retained functionals annihilate its kernel.
Weighted Hilbert duality then gives the integral in
(MA-PR.53).  Absolute convergence of its entries follows
from Cauchy--Schwarz applied to the two continuous jets.
The polar decomposition
\(\mathscr R_n=K_n^{1/2}W_n\) has \(W_nW_n^*=1\).
Its coisometry carries the unit ball onto the Euclidean
unit ball, so
\[
                  \|L\mathscr R_n\|=\|LK_n^{1/2}\|.
 \tag{MA-PR.55}
\]
This remains true for the Banach target \(J_n\).
In particular (MA-PR.53) computes the norm on all of
\(Q_H\), not just the inherited norm on \(V_n\).

The feasible set is a closed affine subspace of the
finite-dimensional space of module maps.
Since \(K_n^{1/2}\) is invertible, the right-hand side
of (MA-PR.55) is a genuine norm on that space.
A minimizing sequence is therefore bounded, has a
convergent subsequence, and its limit is feasible.
This proves attainment and the exact Gram formula.

\emph{Necessity of the global condition.}
Suppose that \(S_g\) is continuous and satisfies
(MA-PR.51).  Put \(F_n=\mathfrak R_nS_g\).
For every \(k\), continuity into (MA-PR.44) gives
\[
             \sup_n\nu_n^k\|F_n\|_{\mathcal L(Q_H,J_n)}
                           <\infty .
 \tag{MA-PR.56}
\]
We must verify that these maps factor through the
actual \(P_n\); orthogonality is not being assumed.
On each finite-dimensional \(V_m\), equivariance
intertwines the matrices \(e^{t(w-a)}\).
Differentiating these finite matrix identities shows
that \(F_n|_{V_m}\) is a module map.
For \(m\ne n\), the two annihilating polynomials have
disjoint primary support, so this map is zero.
For \(m=n\), use the isomorphism in (MA-PR.47)
to write \(F_n|_{V_n}=L_n\mathscr R_n|_{V_n}\).
The algebraic span of the \(V_m\) is dense by
MA-WP-001.  Boundedness of \(F_n\), \(P_n\), and
\(\mathscr R_n\) therefore proves on the whole quotient
\[
                  F_n=F_nP_n=L_n\mathscr R_n .
 \tag{MA-PR.57}
\]
Applying the retained jets to \(iS_g=C_g^{Q_H}\)
and using (MA-PR.47) gives
\(\pi_nL_n=M_{\pi_nb_n(g)}\).
The finite module calculation proves eligibility,
and \(c_n(g)\le\|F_n\|\) proves (MA-PR.52).
This proves necessity without replacing \(P_n\)
by an orthogonal projection.

\emph{Sufficiency and the actual lift.}
Now assume eligibility and (MA-PR.52).
Choose a minimizer \(L_n\) in each block and set
\(F_n=L_n\mathscr R_n=F_nP_n\).
Define
\[
       \widetilde S_gv=\sum_n\mathcal T_nF_nv
              \ \in\mathcal E,\qquad
       S_g=q_\infty\widetilde S_g .
 \tag{MA-PR.58}
\]
The polynomial bounds in (MA-PR.44) and the rapid
decay (MA-PR.52) make this series absolutely convergent
in every defining seminorm, uniformly for \(\|v\|\le1\).
It is a continuous linear lift.
The full jet of its output in block \(n\) is \(F_nv\).
Consequently its retained jets are
\(\pi_nF_nv=M_{\pi_nb_n(g)}\mathscr R_nv\),
which are exactly the retained jets of \(C_g^{Q_H}v\).
The common retained-jet kernel on \(Q_H\) is zero
by (MA-WP.16) and its complete block description.
Hence \(iS_g=C_g^{Q_H}\) as an actual bounded-operator
identity.  The module identities for each \(L_n\)
and the full smooth jet kernel prove the equivariance
in (MA-PR.51).

For the topology of nuclearity, choose an Auerbach
basis \(e_{n,j}\), \(\ell_{n,j}\) of \(J_n\).
Omit blocks for which \(F_n=0\), and put
\[
 \begin{aligned}
 u_{n,j}&=c_n(g)^{1/2}\mathcal T_ne_{n,j},&
 \lambda_{n,j}&=c_n(g)^{1/2},\\
 \alpha_{n,j}(v)&=\ell_{n,j}(F_nv)/c_n(g),&
 \widetilde S_gv&=\sum_{n,j}
                 \lambda_{n,j}u_{n,j}\alpha_{n,j}(v).
 \end{aligned}
 \tag{MA-PR.59}
\]
The functionals have norm at most one.
The target vectors form one bounded set in \(\mathcal E\),
by (MA-PR.44) and rapid decay of \(c_n(g)^{1/2}\).
The coefficients, repeated \(d_n=O(\log\nu_n)\) times,
belong to every \(\ell^p\), \(p>0\).
This proves nuclearity of every positive order with
a common bounded target disk.
For an arbitrary continuous return, use its actual
norms \(\|F_n\|\) from (MA-PR.56) in the same argument.
The resulting lifted series has the same full jets
as that return, so gives that return after \(q_\infty\).
This proves the last assertion for every continuous
solution, not just the chosen minimizers.
\end{proof}

\paragraph{MA-PR-007 (adaptive returns at every width).}
For these fixed Hilbert parameters there are continuous
linear regularizers
\(\mathcal D_N:\mathcal E\to\mathcal E\), chosen
independently of the test, with the following properties.
They preserve \(\mathcal E_{\rm el}\), preserve \(g(0)\),
and satisfy \(\mathcal D_Ng\to g\) in \(\mathcal E\),
uniformly on bounded subsets.
For every eligible \(g\), the test
\(h_{N,g}=\mathcal D_Ng\) has an actual return
\[
 \begin{aligned}
 S_{N,g}:Q_H&\longrightarrow Q_\infty,&
 iS_{N,g}&=C_{h_{N,g}}^{Q_H},\\
 B_{N,g}&=S_{N,g}i,&
 E_{N,g}&=C_{T_{-a}h_{N,g}}^{Q_\infty}-B_{N,g}.
 \end{aligned}
 \tag{MA-PR.60}
\]
Each return has an actual \(\mathcal E\)-valued lift
and is nuclear of every positive order.
For fixed \(N\) these nuclear bounds are uniform for
eligible \(g\) in a bounded subset of \(\mathcal E\).
The construction is linear in \(g\).
The tests \(h_{N,g}\) need not be compactly supported.
No bound uniform in \(N\) for the return maps is claimed.
Eligibility is essential even for approximation: since
\(\mathcal E_{\rm el}\) is closed, no ineligible \(g\)
can be an \(\mathcal E\)-limit of tests each admitting
a continuous equivariant return.  MA-PR-006 makes
every such test eligible.

On the actual Hilbert quotient and on the full smooth
quotient, respectively,
\[
 \begin{aligned}
 C_{h_{N,g}}^{Q_H}&\longrightarrow C_g^{Q_H}
                                      &&\text{in trace norm},\\
 \operatorname{Tr}_{Q_\infty}E_{N,g}
       &\longrightarrow
        \operatorname{Tr}_{Q_\infty}C_{T_{-a}g}
                          -\operatorname{Tr}_{Q_H}C_g .
 \end{aligned}
 \tag{MA-PR.61}
\]
These are operator constructions and trace limits even
when (MA-PR.52) fails for \(g\).
No convergence of \(S_{N,g}\), \(B_{N,g}\), or \(E_{N,g}\)
as maps into \(Q_\infty\) is inferred.

In the rational centered case, for every eligible
compact smooth \(g\) with \(g(0)=0\), the actual
ambient cutoff limit has the transport
\[
 \lim_R\operatorname{Tr}(M_{\chi_R}C_gP_G)
       =A_0(g)+B(g)-P(g)
                       +\lim_{N\to\infty}
                              \operatorname{Tr}_{Q_\infty}E_{N,g}.
 \tag{MA-PR.62}
\]
This now holds at arbitrary widths with the actual
edge multiplicities.  It does not eliminate its correction.

\begin{proof}
\emph{A fixed linear selector in each finite block.}
Let \(\mathcal A_n\subset J_n\) be the subspace satisfying
the local conditions (MA-PR.49).
Choose the following explicit linear solution of the
finite module problem for \(b\in\mathcal A_n\).
At a full retained factor take multiplication by \(b\).
At a proper factor with \(m-r<r\), take
\(c=b\bmod z^r\), represented by its polynomial of
degree less than \(r\), and set \(L(y)=cy\) in \(A\).
Eligibility makes \(c\) divisible by \(z^{m-r}\).
If \(m-r\ge r\), eligibility makes the multiplier
on \(B\) zero and take \(L=0\).
Take zero at discarded factors.
Denote the resulting linear selector by
\(\mathcal L_n(b):B_n\to J_n\).
All choices here are finite-dimensional and satisfy
\[
       \pi_n\mathcal L_n(b)=M_{\pi_nb}.
 \tag{MA-PR.63}
\]
Measure its exact cost by
\[
 A_n=\max\left\{1,\
   \sup_{\substack{b\in\mathcal A_n\\\|b\|_n\le1}}
       \|\mathcal L_n(b)\mathscr R_n\|_{\mathcal L(Q_H,J_n)}
                                    \right\}<\infty .
 \tag{MA-PR.64}
\]
The supremum is finite by finite dimensionality.
It uses the full source norm in (MA-PR.53), hence
retains all oblique-projection and near-edge costs.
No estimate for the growth of \(A_n\) is required.

Set
\[
 \theta_{N,n}=
 \begin{cases}
 1,& |n|\le N,\\
 e^{-\nu_n}/A_n,& |n|>N .
 \end{cases}
 \qquad 0<\theta_{N,n}\le1 .
 \tag{MA-PR.65}
\]
Using the centered lifts of MA-WP, define first
\[
 D_Ng=g-\sum_n(1-\theta_{N,n})
                         T_a\mathcal T_n b_n(g).
 \tag{MA-PR.66}
\]
The series converges absolutely in every
\(\mathcal E\) seminorm by (MA-WP.8)--(MA-WP.10).
Its full block class is
\(\theta_{N,n}b_n(g)\).
Thus \(D_N\) is linear, continuous, and preserves
eligibility.  Its positive scalars do not delete a
nonzero block class.

More explicitly, for every defining seminorm \(p\)
there are \(C,D\), independent of \(N\) and \(g\), with
\[
 p(D_Ng-g)
       \le C\sum_{|n|>N}\nu_n^D\|b_n(g)\|_n
                              \longrightarrow0 .
 \tag{MA-PR.67}
\]
The block restrictions of a bounded subset of
\(\mathcal E\) decrease rapidly with uniform constants,
so the displayed convergence is uniform on such a set.
This estimate concerns the regularizers in
\(\mathcal E\), not the return norms.

\emph{Preserving the value-zero test class without changing jets.}
There is a fixed \(b\in I_a\) with \(b(0)=1\).
Indeed let \(\varphi=T_a\psi\ne0\) and
\(\varphi^\dagger(x)=\overline{\varphi(-x)}\).
The full jet kernel of MT-M18 gives
\(\varphi\in I_a\), and the convolution ideal property gives
\[
          b=\frac{\varphi*\varphi^\dagger}
                         {\int_{\mathbb R}|\varphi(x)|^2\,dx}
                   \in I_a,\qquad b(0)=1 .
 \tag{MA-PR.68}
\]
No stability of the ideal under \(\dagger\) is used:
only the first convolution factor must belong to the ideal.
Define
\[
       \mathcal D_Ng
            =D_Ng-\bigl(D_Ng(0)-g(0)\bigr)b .
 \tag{MA-PR.69}
\]
The correction leaves every full spectral jet of \(D_Ng\)
unchanged and restores the value \(g(0)\).
Continuity of evaluation and (MA-PR.67) prove all the
stated \(\mathcal E\)-convergence properties.
In particular the norm-one class \(g(0)=0\) is preserved.
Neither \(b\) nor the resulting regularizers are asserted
to have compact support.

\emph{The actual nuclear maps.}
For eligible \(g\) set
\[
 F_{N,n,g}=\theta_{N,n}\mathcal L_n(b_n(g))\mathscr R_n,
 \qquad
 \widetilde S_{N,g}
       =\sum_n\mathcal T_nF_{N,n,g}.
 \tag{MA-PR.70}
\]
Every \(F_{N,n,g}=F_{N,n,g}P_n\) and is an equivariant
solution for the full class of \(h_{N,g}\).
Its actual source norm is bounded by
\[
 \|F_{N,n,g}\|
   \le
   \begin{cases}
       A_n\|b_n(g)\|_n,& |n|\le N,\\
       e^{-\nu_n}\|b_n(g)\|_n,& |n|>N .
   \end{cases}
 \tag{MA-PR.71}
\]
For fixed \(N\) this sequence decreases rapidly.
Consequently (MA-PR.58)--(MA-PR.59), now with these
actual norms or their positive upper bounds, prove
the convergence of (MA-PR.70), its all-order nuclearity,
and the actual factorization (MA-PR.60).
Linearity follows from the fixed linear selectors.

For completeness the uniformity in the statement
does not hide a uniform bound for \(A_n\).
There is one continuous seminorm \(p_0\) of
\(\mathcal E\) such that
\(\sup_n\|b_n(g)\|_n\le p_0(g)\), by the Mellin
bound on the fixed strip containing the block contours.
For fixed \(N\), use in (MA-PR.71) the positive scalar
sequence \(A_n\) on its finite head and \(e^{-\nu_n}\)
on its tail.  The square-root nuclear construction
then gives one bounded target disk and coefficients
in every \(\ell^p\), uniformly when \(p_0(g)\) is bounded.
The finite head, and hence these return bounds,
can grow without control as \(N\to\infty\).

\emph{Actual trace cyclicity and the limits.}
Compose the common nuclear series for \(S_{N,g}\)
first with \(i\) on the target and then with \(i\)
on the source.  Its two rank-one trace pairings agree
term by term and are absolutely summable.
The first composition is the actual Hilbert operator
\(C_{h_{N,g}}^{Q_H}\); the second is \(B_{N,g}\).
The approximation property and absolute smooth block
expansion in MT-M18 justify the same nuclear trace
as in MA-HM-002 and MA-PR-003.  Therefore
\[
 \operatorname{Tr}_{Q_H}C_{h_{N,g}}
    =\operatorname{Tr}_{Q_\infty}B_{N,g}
    =\operatorname{Tr}_{Q_\infty}C_{T_{-a}h_{N,g}}
                       -\operatorname{Tr}_{Q_\infty}E_{N,g}.
 \tag{MA-PR.72}
\]
These equalities precede any scalar evaluation of a trace.

MA-WP.14--MA-WP.15 is a continuous map
\(\mathcal E\to\mathcal S_1(Q_H)\).
Together with (MA-PR.67)--(MA-PR.69) it proves
the first limit in (MA-PR.61).
For the full smooth term, the absolute rank-one trace
bound from MT-M19 is
\[
 \left|\operatorname{Tr}_{Q_\infty}
                   C_{T_{-a}(h_{N,g}-g)}\right|
 \le\sum_n d_n\|b_n(h_{N,g}-g)\|_n\longrightarrow0 .
 \tag{MA-PR.73}
\]
Rapid restriction decay and continuity of the transform
justify the convergence.  Subtracting the two limits
in (MA-PR.72) proves the second limit in (MA-PR.61).

For the original compact \(g\) in the rational
specialization, apply the actual ambient theorem
MA-WP-003 to \(g\) itself and Meyer's common-space
identity MA-MEY.15--MA-MEY.16 to its full smooth
operator \(C_{T_{-1/2}g}\).
The first gives the negative of the Hilbert quotient
trace; the second gives \(P-A_0-B\).
Their actual operator identifications and (MA-PR.61)
give (MA-PR.62).
This reasoning asserts no ambient trace-class
extension for noncompact \(h_{N,g}\) and requires
no geometric formula for those regularized tests.

Only after this transport, the finite-dimensional
cyclic identity
\(\operatorname{tr}_{J_n}(L\pi_n)
 =\operatorname{tr}_{B_n}(\pi_nL)\)
evaluates the correction limit as
\[
 \lim_N\operatorname{Tr}_{Q_\infty}E_{N,g}
          =\sum_{\Xi(\rho)=0}
              (m_\rho-r_\rho)G(\rho-a).
 \tag{MA-PR.74}
\]
The local zero count and rapid Mellin decay give
absolute convergence.  At a proper nonzero
boundary truncation, eligibility forces \(G(\rho-a)=0\);
this does not make its actual nilpotent correction
operator zero.  Neither a direct-sum splitting
at that boundary nor a semisimple replacement is used.
\end{proof}

\paragraph{MA-PR-008 (a conditional close-pair topological obstruction).}
Suppose that a block contains a fully retained zero \(\rho\)
and a discarded zero \(\sigma\), so \(r_\rho=m_\rho>0\)
and \(r_\sigma=0\).  Set
\[
 \delta_n(\rho,\sigma)
  =\min\{1,\operatorname{dist}(\rho,\partial\Omega_n),
                  \operatorname{dist}(\sigma,\partial\Omega_n)\}.
 \tag{MA-PR.75}
\]
For every eligible test,
\[
 \begin{aligned}
 c_n(g)&\ge
 \frac{\delta_n(\rho,\sigma)}{2|\rho-\sigma|}
       |G(\rho-a)|\,\|\ell_{\rho-a,0}\|_{H^*},\\
 \ell_{\lambda,0}(h)&=\int e^{\lambda x}h(x)\,dx .
 \end{aligned}
 \tag{MA-PR.76}
\]
In particular, along any sequence of such pairs in blocks
whose heights tend to infinity, an unregularized continuous
return requires
\[
       \frac{|G(\rho_n-a)|}{|\rho_n-\sigma_n|}
                       \quad\hbox{to decrease rapidly in }\nu_n .
 \tag{MA-PR.77}
\]
If that necessary rapid decay fails, no continuous
equivariant return for \(g\) exists.
No such arithmetic pair sequence, or failure of
(MA-PR.77) for an arithmetic function, is asserted here.

\begin{proof}
Take any feasible \(L\) in (MA-PR.50) and let
\(F=L\mathscr R_n\).
At the full retained factor its value coordinate is forced:
\[
                  (Fv)(\rho)=G(\rho-a)\ell_{\rho-a,0}(v).
 \tag{MA-PR.78}
\]
At \(\sigma\) every full output jet is zero.
Indeed the source has no \(\sigma\)-primary factor,
and module maps from every other primary factor
to that factor are zero.
The line segment joining \(\rho\) to \(\sigma\)
lies inside the convex rectangle and has distance
at least \(\delta_n(\rho,\sigma)\) from its boundary.
For any holomorphic representative \(H\) of \(Fv\),
Cauchy's estimate on disks of radius \(\delta_n/2\)
and integration along that segment give
\[
 |H(\rho)-H(\sigma)|
       \le 2|\rho-\sigma|\delta_n^{-1}
                            \sup_{\partial\Omega_n}|H| .
 \tag{MA-PR.79}
\]
Take the infimum over representatives, then the
supremum over the unit ball of the \emph{whole}
\(Q_H\).  The continuous value functional annihilates
\(J_H\), so its quotient norm is exactly its \(H^*\)
norm.  Finally minimize over \(L\).
This proves (MA-PR.76).

The selected-height product bound (M18.7) gives
\(\delta_n(\rho,\sigma)\ge c\nu_n^{-D}\).
To see this directly, every horizontal distance
less than one is a factor in the product in (M18.7);
all factors are at most one.  Each such factor,
even with its multiplicity, is therefore bounded
below by that product.  The vertical edges
at real parts \(-1\) and \(2\) are at least one
from the critical strip.  This proves the asserted
polynomial lower bound for the distance to the
\emph{full} block boundary.  It is not a lower
bound for \(|\rho-\sigma|\).

Furthermore \(\|\ell_{\rho-a,0}\|_{H^*}\ge c'>0\),
uniformly for \(0\le\operatorname{Re}\rho\le1\):
restrict its squared dual integral to \([-1,1]\),
where the fixed weight is bounded above and
the real exponential is bounded below.
If a return exists, (MA-PR.52) and (MA-PR.76)
now imply (MA-PR.77), since multiplication by a
fixed power of \(\nu_n\) preserves rapid decay.
Its failure is thus a precise topological obstruction.
\end{proof}

\paragraph{Scope and placement.}
MA-PR-006 separates the finite algebraic condition
from the global continuity condition.  MA-PR-007
closes the exceptional-width regularized return problem
for every locally eligible test and every \(s\ge0\),
including the requested polynomial weights \(s\ge3\).
The remaining unregularized question is exactly
the actual cost condition (MA-PR.52); local eligibility
alone is not promoted to that bound.
The constructed maps land in the full smooth quotient.
At a proper truncated boundary factor there is no
invariant section on which \(i\) is an isomorphism,
as MA-PR-004 proves.
There is therefore no claim of a canonical retained
smooth direct summand, nor of the projector identities
from the generic-edge construction.

The finite objects belong to the abelian category of
finite-length \(\mathbb C[w]\)-modules.
Locally the exact sequence
\[
 0\longrightarrow z^r\mathbb C[z]/(z^m)
   \longrightarrow\mathbb C[z]/(z^m)
   \longrightarrow\mathbb C[z]/(z^r)
   \longrightarrow0
 \tag{MA-PR.80}
\]
has lengths \(m-r,m,r\).  Its relation in that category's
Grothendieck group records these lengths; it does not
split the sequence or erase its nilpotent structure.
Finite traces and determinants refer to endomorphisms
of the underlying finite-dimensional complex vector spaces,
with determinants formed by their top exterior powers.
The common nuclear series in (MA-PR.59) and (MA-PR.72)
then implements the trace passage on the actual quotients.
This is the inherited nuclear-trace placement of
cote~1, scans~079--080 and~111--112, as already specified
in MA-HM and MA-WP.  No new reading of those scans,
arithmetic positivity theorem, or weight realization
is claimed.

The exact analytic inputs are the full MT-M18
block theorem (M18.1)--(M18.4), its selected-height
bound (M18.7), and its entire interpolation estimates
(M18.10)--(M18.13), in the supplied
\texttt{block\_arithmetic\_synthesis.tex};
MT-M19 (M19.3)--(M19.4), in the supplied
\texttt{block\_trace\_hilbert.tex};
and the actual quotient, projections, kernel and
trace-norm synthesis of MA-WP.4--MA-WP.18.
The complete MT-M18 proof and the relevant MT-M19
nuclear-series proof were reread for this construction.
The Meyer transport uses the same actual fixed-sector
isomorphisms MT-M22/MT-M24 and the same common-space
trace signs MA-MEY.15--MA-MEY.16 as MA-PR-003.

The half-line finite-Paley--Wiener quotient in MA-GW
has its own causal domain and boundary blocks.
It has not been substituted for the two-sided
\(Q_H\) or for its canonical map \(i\) here.
Likewise none of the regularized returns identifies
the unbounded off-diagonal maps in Meyer's original
common-space projection on a new Hilbert domain.
No assumption A, M, generic-edge avoidance, simplicity,
or new outside preprint is used.  The actual arithmetic
deficiency and its possible absence remain unasserted.
